\documentclass[11pt,a4paper]{article}
\usepackage[utf8]{inputenc}
\usepackage[T1]{fontenc}
\usepackage{geometry}
\usepackage{enumitem}
\usepackage{parskip}
\usepackage{hyperref}

\geometry{margin=1in}

\title{\textbf{Wampanoag Planting Phenology: The White Oak Leaf Rule}}
\author{}
\date{}

\begin{document}

\maketitle

\section*{Overview}
The well-known Eastern U.S.\ Indigenous agricultural tradition tracks \textbf{white oak leaves} rather than acorns. Specifically, the \textbf{Wampanoag} (through \textbf{Tisquantum}, or Squanto) taught the Plymouth Pilgrims in 1621 to plant corn \textbf{when the young leaves of the white oak tree were the size of a squirrel's ear} (or mouse's ear).

\textit{Note: Acorns ripen and drop in the autumn rather than appearing in spring, so the oak's spring leaf-out served as the seasonal planting clock.}

\section{The Historical Story (Plymouth, 1621)}
\begin{itemize}[leftmargin=1.5em]
    \item \textbf{The Context:} When English colonists arrived at Plymouth (Patuxet) in winter 1620--1621, they brought European seeds (wheat and peas) that struggled in New England soil and had no experience cultivating native maize (corn).
    \item \textbf{The Teaching:} In the spring of 1621, \textbf{Tisquantum (Squanto)}---a Patuxet Wampanoag man who had survived abduction and epidemic disease that devastated his home village---showed the colonists how to read local \textbf{phenology} (seasonal biological cues) rather than relying on a European calendar.
    \item \textbf{The Rule of Thumb:} He instructed them to watch the white oak trees (\textit{Quercus alba}). Once the emerging spring leaves grew to the size of a \textbf{squirrel's ear}, it was time to plant corn seeds in small mounds alongside herring or alewife fish for fertilizer.
\end{itemize}

\section{Scientific Basis of the Phenological Cue}
Connecting directly to \textbf{phenology and remote sensing} (such as tracking the Start of Season on an NDVI time-series curve), this Traditional Ecological Knowledge (TEK) aligns closely with biophysical thresholds:
\begin{itemize}[leftmargin=1.5em]
    \item \textbf{Soil Temperature Threshold:} White oaks are deep-rooted trees that respond to cumulative spring warmth (growing degree days) and photoperiod rather than brief warm spells. When their leaves reach approximately 1--1.5~inches (squirrel-ear size), soil temperatures at planting depth have reliably reached $55^\circ\text{F}$--$60^\circ\text{F}$ ($13^\circ\text{C}$--$15^\circ\text{C}$), the minimum warmth required for corn germination without seed rot.
    \item \textbf{Frost Safety:} By the time white oaks leaf out to that stage, the risk of a killing late-spring frost in Eastern North America has almost always passed.
    \item \textbf{Enduring Legacy:} Colonial farmers and Appalachian homesteaders adopted the Wampanoag adage verbatim, passing down \textit{``Plant corn when oak leaves are the size of a squirrel's (or mouse's) ear''} for over four centuries.
\end{itemize}

\end{document}
